% In this file you should put the actual content of the blueprint.
% It is used both by the web and the print version.
% It should *not* include the \begin{document}.

\chapter*{Introduction}

This blueprint documents the abstract theory of Hecke rings developed in
\texttt{LeanModularForms}, its specialisation to $\mathrm{GL}_2$, and the application to
Strong Multiplicity One. The narrative runs as follows.

A \emph{Hecke pair} $(G,H,\Delta)$ produces a free $\mathbb{Z}$-module on the double
cosets $H\backslash\Delta/H$ (Chapter~\ref{chap:hecke-pairs}). A convolution product
whose structure constants count coset overlaps makes this an associative unital ring
(Chapter~\ref{chap:ring-structure}), carrying a degree homomorphism to $\mathbb{Z}$
(Chapter~\ref{chap:degree}). An \emph{anti-involution} of the pair that fixes every
double coset forces the ring to be commutative
(Chapter~\ref{chap:commutativity}); for $\mathrm{GL}_n$ the transpose $g\mapsto {}^{t}g$
is such an anti-involution, so the $\mathrm{GL}_n$ Hecke algebra is commutative.
Specialising to $\mathrm{GL}_2$ gives the operators $T(a,d)$ and
$T(m)=\sum_{a\mid m}T(a,m/a)$ (Chapter~\ref{chap:gl2-operators}), which obey Shimura's
multiplication table, Theorem~3.24 (Chapter~\ref{chap:multiplication-table}). Finally,
commutativity makes simultaneous diagonalisation possible, and the rigidity of the
resulting eigenvalue systems yields Strong Multiplicity One, Miyake Theorem~4.6.12
(Chapter~\ref{chap:smo}).

Every declaration referenced below is fully formalised and \texttt{sorry}-free in the
\texttt{LeanModularForms} project; the \texttt{\textbackslash leanok} marks reflect this.

\chapter{Hecke pairs and double cosets}\label{chap:hecke-pairs}

The basic combinatorial datum of Hecke theory is a \emph{Hecke pair} $(G,H,\Delta)$, from
which one builds two free $\mathbb{Z}$-modules: one on the double cosets
$H\backslash\Delta/H$, which will carry the Hecke ring structure, and one on the left
cosets $Hg$, on which that ring will act. This chapter records these definitions together
with the structural fact, following [Sh, \S3.1], that each double coset $HgH$ is a finite
disjoint union of left cosets.

\begin{definition}[Hecke pair]\label{def:hecke-pair}
\lean{HeckeRing.HeckePair}
\leanok
A \emph{Hecke pair} $(G,H,\Delta)$ consists of a group $G$, a subgroup $H$ of $G$, and a
submonoid $\Delta$ of $G$ satisfying $H\le\Delta$ and $\Delta\le\operatorname{Comm}(H)$,
where $\operatorname{Comm}(H)$ is the commensurator of $H$ in $G$. Equivalently, every
$g\in\Delta$ is commensurable with the identity in the sense that
$H\cap gHg^{-1}$ has finite index in both $H$ and $gHg^{-1}$.
\end{definition}

\begin{definition}[Double cosets]\label{def:hecke-coset}
\lean{HeckeRing.HeckeCoset}
\uses{def:hecke-pair}
\leanok
For a Hecke pair $(G,H,\Delta)$, the set of \emph{double cosets} $H\backslash\Delta/H$ is
the quotient of $\Delta$ by the equivalence relation that identifies $g$ and $g'$ whenever
$HgH=Hg'H$, that is, whenever $g'=h_1 g h_2$ for some $h_1,h_2\in H$. We write a typical
double coset as $HgH$.
\end{definition}

\begin{definition}[Left cosets in $\Delta$]\label{def:hecke-left-coset}
\lean{HeckeRing.HeckeLeftCoset}
\uses{def:hecke-pair}
\leanok
For a Hecke pair $(G,H,\Delta)$, the set of \emph{left cosets} is the quotient of $\Delta$
by the equivalence relation that identifies $g$ and $g'$ whenever $Hg=Hg'$. We write a
typical left coset as $Hg$.
\end{definition}

\begin{lemma}[Double coset as a union of left cosets]\label{lem:dc-eq-iUnion-lc}
\lean{HeckeRing.DoubleCoset.doubleCoset_eq_iUnion_leftCosets}
\uses{def:hecke-pair}
\leanok
Let $g\in\Delta$. Then the double coset $HgH$ is the union
\[
  HgH \;=\; \bigcup_{i\,\in\, H/(H\cap gHg^{-1})} H\,(g_i\,g),
\]
where $i$ ranges over the coset space $H/(H\cap gHg^{-1})$ and $g_i\in H$ is a
representative of $i$. In particular $HgH$ is a disjoint union of finitely many left
cosets of $H$.
\end{lemma}

\begin{proof}
\uses{def:hecke-pair}
\leanok
Writing $K=H\cap gHg^{-1}$, the subgroup $H$ is the disjoint union of the left cosets
$g_i K$ as $i$ ranges over $H/K$. Multiplying this decomposition on the right by the set
$gHg^{-1}$ and using that $K$ stabilises $gHg^{-1}$ under left multiplication, one obtains
$H\,gHg^{-1}=\bigcup_i g_i\,gHg^{-1}$; cancelling $g^{-1}$ on the right and translating by
$g$ turns the left-hand side into $HgH$ and each summand into the left coset $H(g_i g)$.
Commensurability of $g$ with the identity, which holds because
$\Delta\le\operatorname{Comm}(H)$, makes the index $[H:H\cap gHg^{-1}]$ finite, so the
union is finite, and distinct cosets $g_iK$ give distinct left cosets $H(g_i g)$, whence
the union is disjoint.
\end{proof}

\begin{definition}[Left-coset index of a double coset]\label{def:decomp-quot}
\lean{HeckeRing.decompQuot}
\uses{def:hecke-pair, lem:dc-eq-iUnion-lc}
\leanok
For $g\in\Delta$, the \emph{decomposition index set} of $g$ is the coset space
$H/(H\cap gHg^{-1})$, which indexes the left cosets appearing in the decomposition
$HgH=\bigsqcup_i H(g_i g)$ of Lemma~\ref{lem:dc-eq-iUnion-lc}. Since
$\Delta\le\operatorname{Comm}(H)$, this index set is finite.
\end{definition}

\begin{definition}[Hecke ring, underlying module]\label{def:hecke-ring-carrier}
\lean{HeckeRing.𝕋}
\uses{def:hecke-coset}
\leanok
Let $(G,H,\Delta)$ be a Hecke pair. The \emph{Hecke ring}
$\mathbb{T}=\mathcal{H}(G,H,\Delta)$ is, as a $\mathbb{Z}$-module, the free
$\mathbb{Z}$-module on the set of double cosets $H\backslash\Delta/H$; that is, the
$\mathbb{Z}$-valued functions on $H\backslash\Delta/H$ of finite support. Its
multiplication is introduced in Chapter~\ref{chap:ring-structure}.
\end{definition}

\begin{definition}[Hecke module]\label{def:hecke-module}
\lean{HeckeRing.HeckeModule}
\uses{def:hecke-left-coset}
\leanok
Let $(G,H,\Delta)$ be a Hecke pair. The \emph{Hecke module} is the free
$\mathbb{Z}$-module on the set of left cosets $Hg$; that is, the $\mathbb{Z}$-valued
functions of finite support on the set of left cosets. This is the module on which the
Hecke ring $\mathbb{T}$ acts.
\end{definition}

\chapter{The Hecke ring structure}\label{chap:ring-structure}

Fix a Hecke pair $(G,H,\Delta)$. This chapter equips the free $\mathbb{Z}$-module
$\mathbb{T}$ on the double cosets $H\backslash\Delta/H$ with its convolution product,
whose structure constants are Shimura's coset-overlap multiplicities, and assembles the
result into a unital associative ring [Sh, \S3.1].

\begin{definition}[Product map on representatives]\label{def:mulMap}
\lean{HeckeRing.mulMap}
\uses{def:decomp-quot, def:hecke-coset}
\leanok
Let $g,g'\in\Delta$, and choose left-coset representatives so that
$HgH=\bigsqcup_i Hg_i$ and $Hg'H=\bigsqcup_j Hg'_j$. The \emph{product map} sends a pair
of representatives $(g_i,g'_j)$ to the double coset $H(g_i\,g'_j)H$ of their product.
\end{definition}

\begin{definition}[Shimura's multiplicity]\label{def:hecke-multiplicity}
\lean{HeckeRing.heckeMultiplicity}
\uses{def:decomp-quot}
\leanok
For $g,g',d\in\Delta$ define the \emph{multiplicity} $\mu(g,g';d)$ to be the number of
pairs of left-coset representatives $(g_i,g'_j)$, drawn from the decompositions
$HgH=\bigsqcup_i Hg_i$ and $Hg'H=\bigsqcup_j Hg'_j$, that satisfy
$g_i\,g'_j\,H=d\,H$. This is the structure constant of the Hecke ring
[Sh, \S3.1, Prop 3.2].
\end{definition}

\begin{definition}[Structure-constant coefficients]\label{def:m-finsupp}
\lean{HeckeRing.m}
\uses{def:hecke-multiplicity}
\leanok
For $g,g'\in\Delta$ let $m(g,g')$ be the finitely-supported function on double cosets
$d\mapsto\mu(g,g';d)$. Only finitely many double cosets occur as $H(g_i\,g'_j)H$, so this
is well defined; it records the coefficients of the product of the two basis elements.
\end{definition}

\begin{definition}[Basis element]\label{def:T-single}
\lean{HeckeRing.T_single}
\uses{def:hecke-ring-carrier}
\leanok
For a double coset $D=HgH$ and a coefficient $b$, let $T_D$ (also written $[D]$) be the
corresponding basis element of $\mathbb{T}$: the function on double cosets taking the
value $b$ at $D$ and $0$ elsewhere. The elements $T_D$ with $b=1$ form a $\mathbb{Z}$-basis
of $\mathbb{T}$.
\end{definition}

\begin{lemma}[Convolution product]\label{lem:mul-def}
\lean{HeckeRing.mul_def}
\uses{def:hecke-multiplicity, def:T-single}
\leanok
For $f,g\in\mathbb{T}$ the product is the convolution
\[
  f\cdot g \;=\; \sum_{D_1}\sum_{D_2} f(D_1)\,g(D_2)\,\bigl(T_{D_1}\cdot T_{D_2}\bigr),
\]
the sums running over the (finite) supports of $f$ and $g$, with each basis product
$T_{D_1}\cdot T_{D_2}$ expanded through the structure constants $\mu(D_1,D_2;\,\cdot\,)$.
\end{lemma}

\begin{proof}
\uses{def:hecke-multiplicity, def:T-single}
\leanok
Both $f$ and $g$ are finite $\mathbb{Z}$-linear combinations of basis elements. The
product is defined to be $\mathbb{Z}$-bilinear, so it distributes over these
combinations, and the displayed double sum is exactly its unfolding on the supports. The
inner term $T_{D_1}\cdot T_{D_2}$ is, by definition, the coefficient function
$m(D_1,D_2)$ collecting the multiplicities $\mu(D_1,D_2;d)$. Hence the identity holds by
definition of the convolution.
\end{proof}

\begin{lemma}[Product of basis elements]\label{lem:T-single-mul}
\lean{HeckeRing.T_single_mul_T_single}
\uses{def:T-single, def:hecke-multiplicity, def:m-finsupp}
\leanok
For double cosets $D_1,D_2$,
\[
  [D_1]\cdot[D_2] \;=\; \sum_{d} \mu(D_1,D_2;d)\,[d],
\]
the sum running over double cosets $d$, with at most finitely many nonzero terms. More
generally, for coefficients $a,b$ one has
$T_{D_1}\cdot T_{D_2}=a\,b\cdot m(D_1,D_2)$.
\end{lemma}

\begin{proof}
\uses{def:T-single, def:hecke-multiplicity, def:m-finsupp}
\leanok
Apply the convolution formula of Lemma~\ref{lem:mul-def} to $f=[D_1]$ and $g=[D_2]$. Each
is supported on a single double coset, so the double sum collapses to the single term
$a\,b$ times the coefficient function $m(D_1,D_2)$. By definition of
Definition~\ref{def:m-finsupp} the value of $m(D_1,D_2)$ at $d$ is the multiplicity
$\mu(D_1,D_2;d)$, giving the displayed expansion.
\end{proof}

\begin{lemma}[Module action as a coset sum]\label{lem:smul-eq-sum}
\lean{HeckeRing.smul_eq_sum}
\uses{def:hecke-module, def:decomp-quot}
\leanok
The Hecke ring $\mathbb{T}$ acts on the left-coset module. For $T\in\mathbb{T}$ and a
formal sum $\nu$ of left cosets, the action is the double sum
\[
  T\cdot\nu \;=\; \sum_{D_1}\sum_{\beta} T(D_1)\,\nu(\beta)
    \sum_{i} \bigl[\,g_i\,\beta\,\bigr],
\]
where $D_1$ ranges over the support of $T$, $\beta$ over the support of $\nu$, and $g_i$
over left-coset representatives of $D_1$, each summand $[\,g_i\,\beta\,]$ being the left
coset obtained by translating $\beta$ by the representative $g_i$.
\end{lemma}

\begin{proof}
\uses{def:hecke-module, def:decomp-quot}
\leanok
This is the unfolding of the definition of the action. The action is $\mathbb{Z}$-bilinear
in $T$ and $\nu$, so it distributes over their basis expansions; on a pair of basis
elements it sends $([D_1],[\beta])$ to the sum of the left cosets $[\,g_i\,\beta\,]$ over
the representatives $g_i$ of $D_1$, which is precisely the displayed orbit sum.
\end{proof}

\begin{lemma}[Associativity]\label{lem:mul-assoc}
\lean{HeckeRing.mul_assoc_𝕋}
\uses{lem:T-single-mul, lem:smul-eq-sum}
\leanok
The product on $\mathbb{T}$ is associative: $(f\cdot g)\cdot h=f\cdot(g\cdot h)$ for all
$f,g,h\in\mathbb{T}$.
\end{lemma}

\begin{proof}
\uses{lem:T-single-mul, lem:smul-eq-sum}
\leanok
Associativity is deduced from the action of $\mathbb{T}$ on the left-coset module of
Lemma~\ref{lem:smul-eq-sum}. That action is a faithful scalar tower: ring multiplication
is compatible with the module action, so $((f\cdot g)\cdot h)\cdot\nu$ and
$(f\cdot(g\cdot h))\cdot\nu$ agree for every formal sum $\nu$ of left cosets. The two
agree because each amounts to acting successively by $h$, then $g$, then $f$, a
re-indexing of the triple coset sum that is manifestly independent of bracketing. By
faithfulness of the action, the two ring elements are equal. Concretely on basis elements
this is the double-coset counting identity behind Lemma~\ref{lem:T-single-mul}.
\end{proof}

\begin{lemma}[Identity element]\label{lem:one-def}
\lean{HeckeRing.one_def}
\uses{def:T-single, def:hecke-coset}
\leanok
The multiplicative identity of $\mathbb{T}$ is $[H]=T_{[H]}$, the basis element attached
to the trivial double coset $H=H\cdot 1\cdot H$ with coefficient $1$.
\end{lemma}

\begin{proof}
\uses{def:T-single, def:hecke-coset}
\leanok
The trivial double coset $H$ has the single left coset $H$ as its decomposition, so for
any double coset $D$ the multiplicities $\mu(H,D;\,\cdot\,)$ and $\mu(D,H;\,\cdot\,)$ are
$1$ at $D$ and $0$ elsewhere. By Lemma~\ref{lem:T-single-mul} this gives
$[H]\cdot[D]=[D]=[D]\cdot[H]$, and bilinearity (Lemma~\ref{lem:mul-def}) extends the
two-sided identity property from basis elements to all of $\mathbb{T}$.
\end{proof}

\begin{theorem}[The Hecke ring]\label{thm:hecke-ring}
\lean{HeckeRing.instRing}
\uses{lem:mul-def, lem:mul-assoc, lem:one-def}
\leanok
With the convolution product and identity above, $\mathbb{T}$ is a unital associative ring
(indeed a $\mathbb{Z}$-algebra) [Sh, \S3.1].
\end{theorem}

\begin{proof}
\uses{lem:mul-def, lem:mul-assoc, lem:one-def}
\leanok
The underlying additive group is the free $\mathbb{Z}$-module on double cosets, hence an
abelian group. The convolution product of Lemma~\ref{lem:mul-def} is
$\mathbb{Z}$-bilinear by construction, so distributivity over addition and compatibility
with integer scaling are immediate. Associativity is Lemma~\ref{lem:mul-assoc}, and
Lemma~\ref{lem:one-def} provides a two-sided multiplicative identity. Together these are
exactly the ring axioms.
\end{proof}

\chapter{The degree homomorphism}\label{chap:degree}

Each double coset splits into finitely many left cosets, and counting them defines a ring
homomorphism from the Hecke ring to $\mathbb{Z}$ that records the size of each operator.

\begin{definition}\label{def:hecke-coset-deg}
\lean{HeckeRing.HeckeCoset_deg}\leanok
\uses{def:hecke-coset, def:decomp-quot}
Let $(G,H,\Delta)$ be a Hecke pair and let $D=HgH$ be a double coset. Writing
$D=\bigsqcup_i Hg_i$ for its decomposition into finitely many left cosets, the
\emph{degree} of $D$ is
\[
  \deg D \;=\; [\,H : H\cap gHg^{-1}\,],
\]
the number of left cosets $Hg_i$ comprising $D$.
\end{definition}

\begin{lemma}\label{lem:smulOrbit-card}
\lean{HeckeRing.smulOrbit_card}\leanok
\uses{def:hecke-coset-deg}
Fix a double coset $D=HgH$ with left-coset representatives $g_i$, and let $H\beta$ be any
left coset with $\beta\in\Delta$. Then the orbit
\[
  \{\,Hg_i\beta : i\,\}
\]
has cardinality equal to $\deg D$.
\end{lemma}

\begin{proof}\leanok
\uses{def:hecke-coset-deg}
The map sending the $i$-th representative to the left coset $H(\beta g_i g)$ is a
surjection onto the orbit by construction, so it suffices to show it is injective. If two
representatives produced the same coset, then after left-cancelling $\beta$ the
corresponding cosets $Hg_i g$ and $Hg_j g$ would meet, hence coincide; this contradicts
the fact that distinct representatives index distinct left cosets in the decomposition of
$D$. Thus the orbit has exactly as many elements as there are left cosets in $D$, namely
$\deg D$ [Sh, \S3.1].
\end{proof}

\begin{definition}\label{def:deg}
\lean{HeckeRing.deg}\leanok
\uses{def:hecke-coset-deg, def:T-single, thm:hecke-ring}
The \emph{degree homomorphism} is the ring homomorphism
\[
  \deg : \mathbb{T} \longrightarrow \mathbb{Z}
\]
determined by sending each basis element $[D]$ to $\deg D$ and extending
$\mathbb{Z}$-linearly; that is, $\deg\bigl(\sum_D a_D\,[D]\bigr)=\sum_D a_D\,\deg D$. It
sends the identity to $1$, is additive, and is multiplicative
[Sh, \S3.1, Prop.~3.3].
\end{definition}

\begin{lemma}\label{lem:deg-mul}
\lean{HeckeRing.deg_mul}\leanok
\uses{def:deg, lem:smulOrbit-card, def:hecke-multiplicity}
For all $f,g$ in the Hecke ring,
\[
  \deg(f\cdot g)=\deg(f)\,\deg(g).
\]
\end{lemma}

\begin{proof}\leanok
\uses{def:deg, lem:smulOrbit-card, def:hecke-multiplicity}
Since $\deg$ is additive, it suffices to verify multiplicativity on a pair of basis
elements $[D_1]$ and $[D_2]$. Expanding the product in the Hecke ring via the structure
constants gives $[D_1]\cdot[D_2]=\sum_d \mu(D_1,D_2;d)\,[d]$, so by definition
\[
  \deg\bigl([D_1]\cdot[D_2]\bigr)=\sum_d \mu(D_1,D_2;d)\,\deg d.
\]
The right-hand side counts, with multiplicity, the left cosets occurring in the product
double coset. Grouping these by the right action of $D_2$ on the left cosets of $D_1$ and
applying the orbit-size identity, each orbit contributes exactly $\deg D_2$ cosets, and
there are $\deg D_1$ orbits; hence the total is $\deg D_1\,\deg D_2$ [Sh, \S3.1].
\end{proof}

\chapter{Commutativity via anti-involutions}\label{chap:commutativity}

A Hecke pair may carry an order-reversing symmetry that fixes each of its double cosets; when
it does, the structure constants are symmetric and the Hecke ring is commutative. Applying this
to the transpose on the arithmetic $\mathrm{GL}_n$ pair shows that the $\mathrm{GL}_n$ Hecke
algebra is commutative.

\begin{definition}\label{def:anti-involution}
\lean{HeckeRing.AntiInvolution}\leanok
\uses{def:hecke-pair}
Let $(G,H,\Delta)$ be a Hecke pair. An \emph{anti-involution} of $(G,H,\Delta)$ is an
anti-homomorphism $g\mapsto\bar g$ of $G$ that is involutive and preserves the pair. Explicitly,
it satisfies
\[
  \bar{\bar g}=g,\qquad \overline{ab}=\bar b\,\bar a,
\]
and maps $H$ into $H$ and $\Delta$ into $\Delta$, that is, $\bar g\in H$ whenever $g\in H$ and
$\bar g\in\Delta$ whenever $g\in\Delta$ [Sh, \S3.1].
\end{definition}

\begin{definition}\label{def:onHeckeCoset}
\lean{HeckeRing.AntiInvolution.onHeckeCoset}\leanok
\uses{def:anti-involution, def:hecke-coset}
Let $g\mapsto\bar g$ be an anti-involution of the Hecke pair $(G,H,\Delta)$. It induces a map on
the set $H\backslash\Delta/H$ of double cosets,
\[
  HgH \;\longmapsto\; H\bar gH .
\]
This is well defined: since the anti-involution reverses products and preserves $H$, two
representatives of the same double coset have barred images lying in a single double coset, so
the assignment depends only on $HgH$. The induced map is again involutive [Sh, \S3.1].
\end{definition}

\begin{theorem}\label{thm:mul-comm-of-ai}
\lean{HeckeRing.AntiInvolution.mul_comm_of_antiInvolution}\leanok
\uses{def:onHeckeCoset, def:hecke-multiplicity, thm:hecke-ring}
Let $g\mapsto\bar g$ be an anti-involution of the Hecke pair $(G,H,\Delta)$ whose induced map on
double cosets fixes every double coset, that is, $H\bar gH=HgH$ for all $g\in\Delta$. Then the
Hecke ring $\mathbb{T}$ is commutative: for all $f,g\in\mathbb{T}$,
\[
  f\cdot g \;=\; g\cdot f .
\]
\end{theorem}

\begin{proof}\leanok
\uses{def:onHeckeCoset, def:hecke-multiplicity, thm:hecke-ring}
Since the product on $\mathbb{T}$ is $\mathbb{Z}$-bilinear, it suffices to prove the identity on
a pair of basis elements $[D_1]$ and $[D_2]$, where the product is governed by the structure
constants $\mu(D_1,D_2;d)$. The claim therefore reduces to the symmetry
\[
  \mu(D_1,D_2;d)\;=\;\mu(D_2,D_1;d)
\]
for every double coset $d$. To see this, apply the anti-involution to a left-coset
decomposition $D_1D_2=\bigsqcup_i Hg_i$. Because $g\mapsto\bar g$ reverses products, it carries
this decomposition to one of $D_2D_1$; because it fixes each double coset, the double coset of
each $g_i$ is preserved. Tracking the left coset of a fixed representative $d$ under this
correspondence gives a bijection between the pairs of representatives counted by
$\mu(D_1,D_2;d)$ and those counted by $\mu(D_2,D_1;d)$, so the two multiplicities agree. Hence
the products $[D_1]\cdot[D_2]$ and $[D_2]\cdot[D_1]$ coincide on every basis pair, and
$\mathbb{T}$ is commutative [Sh, \S3.1, Prop.~3.8].
\end{proof}

\begin{definition}\label{def:commring-of-ai}
\lean{HeckeRing.instCommRing_of_antiInvolution}\leanok
\uses{thm:mul-comm-of-ai, thm:hecke-ring}
Given an anti-involution of the Hecke pair $(G,H,\Delta)$ whose induced map fixes every double
coset, the resulting commutativity of multiplication upgrades the ring structure on $\mathbb{T}$
to a \emph{commutative ring} structure, obtained from the ring structure on $\mathbb{T}$ by
supplying the symmetry $f\cdot g=g\cdot f$ from Theorem~\ref{thm:mul-comm-of-ai}.
\end{definition}

\begin{definition}\label{def:gl-pair-ai}
\lean{HeckeRing.GLn.GL_pair_antiInvolution}\leanok
\uses{def:anti-involution, def:GL-pair}
For the arithmetic $\mathrm{GL}_n$ Hecke pair, the \emph{transpose} $g\mapsto{}^{t}g$ is an
anti-involution. It reverses products, since ${}^{t}(ab)={}^{t}b\,{}^{t}a$, and is involutive,
since ${}^{t}({}^{t}g)=g$. Moreover it preserves the pair: transposing a matrix in
$\mathrm{SL}_n(\mathbb{Z})$ stays in $\mathrm{SL}_n(\mathbb{Z})$, and transposing a
positive-determinant integral matrix preserves both integrality of the entries and the value of
the determinant, so the positive-determinant integer monoid is preserved as well [Sh, \S3.1].
\end{definition}

\begin{lemma}\label{lem:gl-pair-fix}
\lean{HeckeRing.GLn.GL_pair_onHeckeCoset_eq}\leanok
\uses{def:gl-pair-ai, def:onHeckeCoset}
The transpose fixes every double coset of the arithmetic $\mathrm{GL}_n$ pair: for all $g$,
\[
  H\,{}^{t}gH \;=\; HgH .
\]
\end{lemma}

\begin{proof}\leanok
\uses{def:gl-pair-ai, def:onHeckeCoset}
Every double coset has a diagonal representative, given by the elementary divisors of $g$ via
the Smith normal form. A diagonal matrix is equal to its own transpose, so the transpose fixes
this representative and therefore fixes the whole double coset. Equivalently, a matrix and its
transpose share the same elementary divisors, hence lie in the same double coset [Sh, \S3.1].
\end{proof}

\begin{theorem}\label{thm:commring-hecke-algebra}
\lean{HeckeRing.GLn.instCommRing_HeckeAlgebra}\leanok
\uses{def:commring-of-ai, def:gl-pair-ai, lem:gl-pair-fix, def:hecke-algebra}
The $\mathrm{GL}_n$ Hecke algebra $\mathcal{H}_n$ is commutative.
\end{theorem}

\begin{proof}\leanok
\uses{def:commring-of-ai, def:gl-pair-ai, lem:gl-pair-fix, def:hecke-algebra}
The transpose is an anti-involution of the arithmetic $\mathrm{GL}_n$ pair
(Definition~\ref{def:gl-pair-ai}) whose induced map fixes every double coset
(Lemma~\ref{lem:gl-pair-fix}). Applying the commutativity criterion of
Definition~\ref{def:commring-of-ai} to this anti-involution endows $\mathcal{H}_n$ with a
commutative ring structure [Sh, \S3.1, Prop.~3.8].
\end{proof}

\chapter{The $\mathrm{GL}_2$ Hecke operators}\label{chap:gl2-operators}

Specialising the abstract construction to the arithmetic pair attached to $\mathrm{GL}_n$
produces the classical Hecke operators. This chapter sets up that pair, its Hecke ring, the
diagonal double cosets that turn out to span it, and, for $n=2$, the generators $T(a,d)$ and
$T(m)$ of Shimura's multiplication table [Sh, \S3.1].

\begin{definition}[The arithmetic Hecke pair]\label{def:GL-pair}
\lean{HeckeRing.GLn.GL_pair}
\uses{def:hecke-pair}
\leanok
For $n\ge 1$ the \emph{arithmetic Hecke pair} for $\mathrm{GL}_n$ is the triple with group
$G=\mathrm{GL}_n(\mathbb{Q})$, subgroup $H=SL_n(\mathbb{Z})$, and submonoid $\Delta$ equal
to the monoid of integer matrices with positive determinant, regarded inside
$\mathrm{GL}_n(\mathbb{Q})$. One has $SL_n(\mathbb{Z})\le\Delta$, since elements of
$SL_n(\mathbb{Z})$ are integral with determinant $1$, and $\Delta\le\operatorname{Comm}(SL_n(\mathbb{Z}))$,
the commensurator of $SL_n(\mathbb{Z})$: an integer matrix $\alpha$ of determinant $d>0$
conjugates the principal congruence subgroup of level $d$ into $SL_n(\mathbb{Z})$, so
$SL_n(\mathbb{Z})$ and $\alpha\,SL_n(\mathbb{Z})\,\alpha^{-1}$ are commensurable
[Sh, \S3.1, Lemma~3.10]. These two inclusions make the triple a Hecke pair.
\end{definition}

\begin{definition}[The $\mathrm{GL}_n$ Hecke algebra]\label{def:hecke-algebra}
\lean{HeckeRing.GLn.HeckeAlgebra}
\uses{def:GL-pair, def:hecke-ring-carrier}
\leanok
For $n\ge 1$ the \emph{$\mathrm{GL}_n$ Hecke algebra} $\mathcal{H}_n$ is the Hecke ring
$\mathbb{T}$ of the arithmetic Hecke pair of $\mathrm{GL}_n$, that is, the free
$\mathbb{Z}$-module on the double cosets $SL_n(\mathbb{Z})\backslash\Delta/SL_n(\mathbb{Z})$
equipped with its convolution product. For $n=2$ this is the ring whose structure is
described by Shimura's Theorem~3.24.
\end{definition}

\begin{definition}[Diagonal double coset]\label{def:T-diag}
\lean{HeckeRing.GLn.T_diag}
\uses{def:GL-pair, def:hecke-coset}
\leanok
Let $a=(a_1,\dots,a_n)$ be an $n$-tuple of positive integers. The associated
\emph{diagonal double coset} is the double coset
\[
  SL_n(\mathbb{Z})\,\operatorname{diag}(a_1,\dots,a_n)\,SL_n(\mathbb{Z})
\]
of the diagonal matrix $\operatorname{diag}(a_1,\dots,a_n)$ in
$SL_n(\mathbb{Z})\backslash\Delta/SL_n(\mathbb{Z})$. Every double coset of the arithmetic
pair admits such a diagonal representative, with elementary divisors $a_1\mid a_2\mid\dots\mid a_n$,
by the theory of Smith normal form [Sh, \S3.1].
\end{definition}

\begin{definition}[Diagonal basis element]\label{def:T-elem}
\lean{HeckeRing.GLn.T_elem}
\uses{def:T-diag, def:T-single, def:hecke-algebra}
\leanok
For an $n$-tuple $a=(a_1,\dots,a_n)$ of positive integers, the \emph{diagonal basis element}
is the standard basis vector of $\mathcal{H}_n$ attached to the diagonal double coset of
$\operatorname{diag}(a_1,\dots,a_n)$; that is, the basis element with coefficient $1$ at that
double coset and $0$ elsewhere.
\end{definition}

\begin{definition}[$T(a,d)$]\label{def:T-ad}
\lean{HeckeRing.GL2.T_ad}
\uses{def:T-elem}
\leanok
Let $n=2$. For positive integers $a,d$ with $a\mid d$, the operator $T(a,d)$ is the diagonal
basis element of $\mathcal{H}_2$ attached to $\operatorname{diag}(a,d)$. For all other pairs
$(a,d)$ — when $a$ or $d$ vanishes, or when $a\nmid d$ — one sets $T(a,d)=0$; the
divisibility constraint $a\mid d$ records the elementary-divisor normalisation of the
diagonal double coset.
\end{definition}

\begin{definition}[$T(p,p)$]\label{def:T-pp}
\lean{HeckeRing.GL2.T_pp}
\uses{def:T-ad}
\leanok
For a positive integer $p$, the operator $T(p,p)$ is defined as $T(p,p)=T(a,d)$ with $a=d=p$.
When $p$ is prime this is the central, scalar double coset of $pI$, the basis element attached
to the scalar matrix $\operatorname{diag}(p,p)$.
\end{definition}

\begin{definition}[$T(m)$]\label{def:T-sum}
\lean{HeckeRing.GL2.T_sum}
\uses{def:T-ad}
\leanok
For a positive integer $m$, Shimura's operator $T(m)$ is the sum
\[
  T(m) \;=\; \sum_{a\,\mid\, m} T\!\left(a,\tfrac{m}{a}\right)
\]
over all positive divisors $a$ of $m$, where each summand is the operator
$T\!\left(a,\tfrac{m}{a}\right)$ of $\mathcal{H}_2$. Only the divisor pairs $(a,m/a)$ with
$a\mid m/a$ contribute, the remaining terms vanishing by definition of $T(a,d)$.
\end{definition}

\begin{lemma}[$T(1,1)$ is the identity]\label{lem:T-ad-one-one}
\lean{HeckeRing.GL2.T_ad_one_one}
\uses{def:T-ad, lem:one-def}
\leanok
In $\mathcal{H}_2$ one has $T(1,1)=1$, the multiplicative identity of the Hecke algebra.
\end{lemma}

\begin{proof}
\uses{def:T-ad, lem:one-def}
\leanok
Since $1\mid 1$ and both entries are positive, $T(1,1)$ is the diagonal basis element of
$\operatorname{diag}(1,1)$, which is the identity matrix. Its diagonal double coset is
therefore the trivial double coset $SL_2(\mathbb{Z})$, and the basis element attached to the
trivial double coset is exactly the identity of the Hecke ring.
\end{proof}

\begin{lemma}[$T(p)$ at a prime]\label{lem:T-sum-prime}
\lean{HeckeRing.GL2.T_sum_prime}
\uses{def:T-sum, def:T-ad}
\leanok
For a prime $p$ one has $T(p)=T(1,p)$.
\end{lemma}

\begin{proof}
\uses{def:T-sum, def:T-ad}
\leanok
The positive divisors of a prime $p$ are $1$ and $p$, so the defining sum for $T(p)$ has the
two terms $T(1,p)$ and $T(p,1)$. The second vanishes: $p\nmid 1$ since $p>1$, so $T(p,1)=0$
by definition. Hence $T(p)=T(1,p)$.
\end{proof}

\chapter{The multiplication table (Shimura 3.24)}\label{chap:multiplication-table}

Because the $\mathrm{GL}_2$ Hecke algebra $\mathcal{H}_2$ is commutative
(Theorem~\ref{thm:commring-hecke-algebra}), its structure is completely pinned down by how the
generators $T(a,d)$ multiply. This chapter formalises Shimura's Theorem~3.24 [Sh, Thm~3.24]: the
identity $T(1)=1$, the prime-power telescoping relation $T(1,p^k)=T(p^k)-T(p,p)\,T(p^{k-2})$, the
recurrence and product formula for the prime-power operators $T(p^r)$, full multiplicativity
$T(m)\,T(n)=\sum_{d\mid(m,n)} d\,T(d,d)\,T(mn/d^2)$, and the degree formulas
$\deg T(p^k)=1+p+\dots+p^k$ and $\deg T(m)=\sigma_1(m)$. The arguments are arithmetic and
combinatorial: telescoping sums of diagonal classes, the central role of the scalar class
$T(p,p)$, coset-counting for the structure constants, and the Chinese Remainder Theorem for the
coprime factorisation.

\begin{lemma}[$T(1)=1$]\label{lem:T-sum-one}
\lean{HeckeRing.GL2.T_sum_one}
\uses{def:T-sum, def:T-ad}
\leanok
In $\mathcal{H}_2$ one has $T(1)=1$, the multiplicative identity of the Hecke algebra.
\end{lemma}

\begin{proof}
\uses{def:T-sum, def:T-ad, lem:T-ad-one-one}
\leanok
The only positive divisor of $1$ is $1$ itself, so the defining sum
$T(1)=\sum_{a\mid 1} T(a,1/a)$ collapses to the single term $T(1,1)$. By
Lemma~\ref{lem:T-ad-one-one} this equals the identity of $\mathcal{H}_2$.
\end{proof}

\begin{theorem}[Shimura 3.24(2): telescoping relation]\label{thm:T-ad-one-ppow-eq}
\lean{HeckeRing.GL2.T_ad_one_ppow_eq}
\uses{def:T-sum, def:T-ad, def:T-pp}
\leanok
Let $p$ be a prime and $k\ge 2$. Then in $\mathcal{H}_2$
\[
  T(1,p^k) \;=\; T(p^k) \;-\; T(p,p)\,T(p^{k-2}).
\]
\end{theorem}

\begin{proof}
\uses{def:T-sum, def:T-ad, def:T-pp}
\leanok
Expand $T(p^k)=\sum_{i=0}^{\lfloor k/2\rfloor} T(p^i,p^{k-i})$ over the diagonal classes whose
elementary divisors are powers of $p$. Multiplying the scalar class $T(p,p)$ into a diagonal
class shifts both elementary divisors up by one, $T(p,p)\,T(p^j,p^{d})=T(p^{j+1},p^{d+1})$, so
$T(p,p)\,T(p^{k-2})=\sum_{j} T(p^{j+1},p^{k-1-j})$, which is exactly the expansion of $T(p^k)$
with the leading $i=0$ term $T(1,p^k)$ removed. Subtracting the two expansions telescopes to the
single surviving term $T(1,p^k)$, giving the stated identity. (Equivalently, the displayed sum is
proved by reindexing the shifted sum via $i=j+1$ and peeling off the bottom term.)
\end{proof}

\begin{theorem}[Key prime computation]\label{thm:T-sum-prime-mul-T-ad}
\lean{HeckeRing.GL2.T_sum_prime_mul_T_ad}
\uses{lem:T-sum-prime, def:T-ad, def:T-pp, def:hecke-multiplicity}
\leanok
Let $p$ be a prime and $k\ge 1$. Then
\[
  T(p)\,T(1,p^k) \;=\; T(1,p^{k+1}) \;+\; c\,T(p,p^k),
  \qquad
  c \;=\;
  \begin{cases}
    p+1, & k=1,\\
    p,   & k\ge 2.
  \end{cases}
\]
\end{theorem}

\begin{proof}
\uses{lem:T-sum-prime, def:T-ad, def:T-pp, def:hecke-multiplicity}
\leanok
Since $T(p)=T(1,p)$ (Lemma~\ref{lem:T-sum-prime}), the product is the convolution of the basis
elements attached to $\operatorname{diag}(1,p)$ and $\operatorname{diag}(1,p^k)$. By the product
formula for the Hecke ring, its coefficient at a double coset $A$ is the structure constant
$\mu$ counting pairs of left-coset representatives whose product lands in a left coset of $A$.
A determinant computation forces every such $A$ to have total elementary divisor product
$p^{k+1}$, and a divisibility analysis of the first elementary divisor (using that an integral
matrix of determinant $1$ has coprime first column) shows that the only double cosets that can
occur are $\operatorname{diag}(1,p^{k+1})$ and $\operatorname{diag}(p,p^k)$. Counting the
overlapping left cosets gives multiplicity $1$ at $\operatorname{diag}(1,p^{k+1})$ and
multiplicity $c$ at $\operatorname{diag}(p,p^k)$, where the count splits according to whether
$k=1$ (the two contributing cosets coincide, giving $p+1$) or $k\ge 2$ (giving $p$). All other
coefficients vanish, which is the claimed identity.
\end{proof}

\begin{theorem}[Prime-power recurrence]\label{thm:T-sum-ppow-recurrence}
\lean{HeckeRing.GL2.T_sum_ppow_recurrence}
\uses{thm:T-ad-one-ppow-eq, thm:T-sum-prime-mul-T-ad, lem:T-sum-prime, def:T-pp}
\leanok
Let $p$ be a prime and $k\ge 1$. Then
\[
  T(p^{k+1}) \;=\; T(p)\,T(p^k) \;-\; p\,\bigl(T(p,p)\,T(p^{k-1})\bigr).
\]
\end{theorem}

\begin{proof}
\uses{thm:T-ad-one-ppow-eq, thm:T-sum-prime-mul-T-ad, lem:T-sum-prime, def:T-pp}
\leanok
Apply the key computation Theorem~\ref{thm:T-sum-prime-mul-T-ad} at exponent $k$ and rewrite the
two diagonal classes through the telescoping relation Theorem~\ref{thm:T-ad-one-ppow-eq}, using
$T(p,p^k)=T(p,p)\,T(1,p^{k-1})$ to express both $T(1,p^k)$, $T(1,p^{k+1})$ and $T(p,p^k)$ in
terms of the $T(p^\bullet)$. For $k=1$ the coefficient $c=p+1$ contributes an extra $+T(p,p)$
which combines with the base values $T(p^0)=1$ and $T(1,p)=T(p)$; for $k\ge 2$ one proceeds by
strong induction, substituting the recurrence already known at the smaller index and using that
$T(p)$ and $T(p,p)$ commute. Collecting the terms and isolating $T(p^{k+1})$ yields the displayed
three-term recurrence.
\end{proof}

\begin{theorem}[Shimura 3.24(3): prime-power product]\label{thm:T-sum-ppow-mul}
\lean{HeckeRing.GL2.T_sum_ppow_mul}
\uses{thm:T-sum-ppow-recurrence, def:T-pp, def:T-sum}
\leanok
Let $p$ be a prime and let $r\le s$. Then
\[
  T(p^r)\,T(p^s) \;=\; \sum_{i=0}^{r} p^i\,T(p^i,p^i)\,T(p^{r+s-2i}).
\]
\end{theorem}

\begin{proof}
\uses{thm:T-sum-ppow-recurrence, def:T-pp, def:T-sum}
\leanok
Strong induction on $r$. For $r=0$ the sum has the single term $T(p^s)$, and for $r=1$ the
statement is precisely the recurrence Theorem~\ref{thm:T-sum-ppow-recurrence} solved for
$T(p)\,T(p^s)$. For the inductive step write $T(p^{r+2})=T(p)\,T(p^{r+1})-p\,T(p,p)\,T(p^r)$ from
the recurrence and apply the induction hypothesis to $T(p^{r+1})\,T(p^s)$ and $T(p^r)\,T(p^s)$.
Each $T(p)\,T(p^{r+1+s-2i})$ splits, again by the recurrence, into a top term
$T(p^{r+2+s-2i})$ and a shifted term $p\,T(p,p)\,T(p^{r+s-2i})$; the shifted terms cancel against
the contribution of $-p\,T(p,p)\,T(p^r)$ except at the two endpoints, where reindexing
$T(p,p)^{i}\cdot T(p,p)=T(p^{i+1},p^{i+1})$ reassembles the sum over $i=0,\dots,r+2$. This is the
asserted product formula.
\end{proof}

\begin{theorem}[Coprime multiplicativity]\label{thm:T-sum-mul-coprime}
\lean{HeckeRing.GL2.T_sum_mul_coprime}
\uses{def:T-sum, def:T-ad, thm:commring-hecke-algebra}
\leanok
Let $m,n$ be positive integers with $(m,n)=1$. Then $T(m)\,T(n)=T(mn)$.
\end{theorem}

\begin{proof}
\uses{def:T-sum, def:T-ad, thm:commring-hecke-algebra}
\leanok
Expand both factors as sums over divisor pairs and multiply out:
$T(m)\,T(n)=\sum_{a\mid m}\sum_{b\mid n} T(a,m/a)\,T(b,n/b)$. Because $(m,n)=1$, the elementary
divisors of the two diagonal classes are coprime, so the corresponding diagonal double cosets
split as a single double coset under the Chinese Remainder Theorem, giving
$T(a,m/a)\,T(b,n/b)=T(ab,(m/a)(n/b))=T(ab,mn/(ab))$. The map $(a,b)\mapsto ab$ is a bijection from
divisor pairs of $(m,n)$ onto divisors of $mn$ (again by coprimality), so the double sum
reindexes to $\sum_{c\mid mn} T(c,mn/c)=T(mn)$. The terms where a divisibility constraint fails
contribute $0$ on both sides, matching up.
\end{proof}

\begin{theorem}[Shimura 3.24: general product]\label{thm:T-sum-mul}
\lean{HeckeRing.GL2.T_sum_mul}
\uses{thm:T-sum-ppow-mul, thm:T-sum-mul-coprime, def:T-pp, def:T-sum}
\leanok
For positive integers $m,n$,
\[
  T(m)\,T(n) \;=\; \sum_{d\,\mid\,(m,n)} d\,T(d,d)\,T\!\left(\frac{mn}{d^2}\right),
\]
the sum being over the positive divisors $d$ of the greatest common divisor $(m,n)$.
\end{theorem}

\begin{proof}
\uses{thm:T-sum-ppow-mul, thm:T-sum-mul-coprime, def:T-pp, def:T-sum}
\leanok
Both sides are multiplicative in the coprime factorisation of $m$ and $n$, so by induction on the
gcd it suffices to treat one prime at a time. Factor out a prime $q\mid(m,n)$, writing
$m=q^a m'$ and $n=q^b n'$ with $q\nmid m'n'$. The coprime multiplicativity
Theorem~\ref{thm:T-sum-mul-coprime} reduces $T(m)\,T(n)$ to $T(q^a)\,T(q^b)$ times the
$q$-free part $T(m')\,T(n')$, to which the induction hypothesis applies. The prime-power factor is
evaluated by Theorem~\ref{thm:T-sum-ppow-mul}, whose terms $q^i\,T(q^i,q^i)\,T(q^{a+b-2i})$ are
exactly the $d=q^i$ contributions to the divisor sum. Recombining the prime-power sum with the
$q$-free divisor sum, using that divisors of $(m,n)$ factor uniquely as $q^i$ times a divisor of
$(m',n')$ and that $T(d,d)$ is multiplicative on coprime $d$, reproduces the single sum over all
$d\mid(m,n)$.
\end{proof}

\begin{theorem}[Shimura 3.24(7): degree of $T(p^k)$]\label{thm:deg-T-sum-prime-pow}
\lean{HeckeRing.GL2.deg_T_sum_prime_pow}
\uses{def:deg, def:T-sum, def:T-ad}
\leanok
Let $p$ be a prime and $k\ge 0$. Then
\[
  \deg T(p^k) \;=\; \sum_{j=0}^{k} p^{j} \;=\; 1+p+p^2+\dots+p^k.
\]
\end{theorem}

\begin{proof}
\uses{def:deg, def:T-sum, def:T-ad}
\leanok
The degree homomorphism sends a diagonal class to its number of left cosets: for $2i<k$ one has
$\deg T(p^i,p^{k-i})=p^{k-2i-1}(p+1)$, while the scalar classes $T(p^i,p^i)$ have degree $1$.
Applying $\deg$ to the expansion $T(p^k)=\sum_{i} T(p^i,p^{k-i})$ and proceeding by strong
induction: for $k\ge 2$ split off the $i=0$ term, contributing $p^{k-1}(p+1)$, and observe that
each remaining term $T(p^{i+1},p^{k-1-i})$ has the same diagonal ratio — hence the same degree —
as the corresponding term of the $T(p^{k-2})$ expansion. Thus the tail equals
$\deg T(p^{k-2})=1+p+\dots+p^{k-2}$ by the induction hypothesis, and
$p^{k-1}(p+1)+(1+\dots+p^{k-2})=1+p+\dots+p^k$, the geometric sum.
\end{proof}

\begin{theorem}[Shimura 3.24(6): degree of $T(m)$]\label{thm:deg-T-sum}
\lean{HeckeRing.GL2.deg_T_sum}
\uses{def:deg, thm:deg-T-sum-prime-pow, thm:T-sum-mul-coprime, lem:deg-mul}
\leanok
For a positive integer $m$,
\[
  \deg T(m) \;=\; \sum_{d\,\mid\, m} d \;=\; \sigma_1(m),
\]
the sum of the positive divisors of $m$.
\end{theorem}

\begin{proof}
\uses{def:deg, thm:deg-T-sum-prime-pow, thm:T-sum-mul-coprime, lem:deg-mul}
\leanok
Both $\deg T(\cdot)$ and $\sigma_1$ are multiplicative arithmetic functions, so it suffices to
match them on prime powers and to use multiplicativity on coprime factors. On a prime power
Theorem~\ref{thm:deg-T-sum-prime-pow} gives $\deg T(p^k)=1+p+\dots+p^k=\sum_{d\mid p^k} d
=\sigma_1(p^k)$. For coprime $a,b$ the coprime multiplicativity $T(ab)=T(a)\,T(b)$
(Theorem~\ref{thm:T-sum-mul-coprime}) together with the ring-homomorphism property
$\deg(fg)=\deg f\,\deg g$ (Lemma~\ref{lem:deg-mul}) yields
$\deg T(ab)=\deg T(a)\,\deg T(b)=\sigma_1(a)\,\sigma_1(b)=\sigma_1(ab)$. Induction on the prime
factorisation of $m$ then gives $\deg T(m)=\sigma_1(m)$.
\end{proof}

\chapter{Strong Multiplicity One}\label{chap:smo}

This is the headline chapter. Strong Multiplicity One [Miy, Thm~4.6.12] (equivalently [DS, Thm~5.8.2])
says that a Hecke eigenform is rigidly determined by its eigenvalues at almost all primes: if a
normalised newform $f$ in $S_k(N,\chi)$ and any cusp-form Hecke eigenfunction $g$ in $S_k(N,\chi)$
share the eigenvalues $\lambda_n$ for all $n$ with $(n,N)=1$ outside a finite set, then $g=c\,f$ for a
single constant $c\in\mathbb{C}$. The result rests squarely on the commutativity of the $\mathrm{GL}_2$
Hecke algebra established in the previous chapters (Theorem~\ref{thm:commring-hecke-algebra}) together
with its multiplication table (Theorem~\ref{thm:T-sum-mul}): commutativity is what makes simultaneous
diagonalisation possible, and the rigidity of the eigenvalue systems is read off from the structure of
the $T(n)$. The proof decomposes $g$ along the Petersson-orthogonal new/old splitting, shows the new
part is a multiple of $f$, and shows the old part vanishes.

\section{Eigenforms, newforms, and the new/old decomposition}

\begin{definition}[Nebentypus space $S_k(N,\chi)$]\label{def:modFormCharSpace}
\lean{HeckeRing.GL2.modFormCharSpace}
\leanok
For a positive integer $N$, a weight $k$, and a Dirichlet character $\chi$ modulo $N$, the
\emph{Nebentypus space} (modular-form version $M_k(N,\chi)$, cusp version $S_k(N,\chi)$) is the
common eigenspace
\[
  M_k(N,\chi) \;=\; \bigcap_{d} \ker\bigl(\langle d\rangle - \chi(d)\bigr),
\]
the intersection over all units $d$ modulo $N$ of the $\chi(d)$-eigenspaces of the diamond operators
$\langle d\rangle$. Equivalently, $f$ lies in $M_k(N,\chi)$ if and only if $\langle d\rangle f=\chi(d)\,f$
for every $d$ coprime to $N$.
\end{definition}

\begin{definition}[Eigenform]\label{def:eigenform}
\lean{HeckeRing.GL2.Eigenform}
\uses{def:modFormCharSpace, thm:commring-hecke-algebra, thm:T-sum-mul}
\leanok
An \emph{eigenform} of level $\Gamma_1(N)$ and weight $k$ is a cusp form $f$ carrying a Nebentypus
character $\chi$, so that $f$ lies in the eigenspace $S_k(N,\chi)$, and which is a simultaneous
eigenfunction of all the Hecke operators $T(n)$ with $(n,N)=1$. Because the Hecke operators commute
(Theorem~\ref{thm:commring-hecke-algebra}, whose product law is the multiplication table
Theorem~\ref{thm:T-sum-mul}), such common eigenfunctions exist in abundance: a finite-dimensional space
of cusp forms is simultaneously diagonalised by the commuting family $\{T(n)\}_{(n,N)=1}$. The eigenvalue
is recorded as $T(n)f=\lambda_n\,f$, where $\lambda_n=\chi(n)\,\nu_n$ relates the classical Hecke
eigenvalue $\lambda_n$ to the underlying Hecke-ring eigenvalue $\nu_n$ through the diamond factor $\chi(n)$.
\end{definition}

\begin{definition}[Old subspace $S_k^{\flat}(N)$]\label{def:cuspFormsOld}
\lean{HeckeRing.GL2.cuspFormsOld}
\leanok
The \emph{old subspace} $S_k^{\flat}(N)$ is the span of all level-raising images: for each proper divisor
$M\mid N$ (with $M\neq N$) and each divisor $\ell$ of $N/M$ with $\ell>1$, the degeneracy map
$f\mapsto f(\ell z)$ carries the cusp space at level $M$ into level $N$, and $S_k^{\flat}(N)$ is the
subspace generated by all such images,
\[
  S_k^{\flat}(N) \;=\; \operatorname{span}\bigl\{\, f(\ell z) \ :\ f\in S_k(\Gamma_1(M)),\ M\mid N,\ M\neq N \,\bigr\}.
\]
This is the [DS, (5.18)] definition of the forms ``coming from lower level''.
\end{definition}

\begin{definition}[New subspace $S_k^{\sharp}(N)$]\label{def:cuspFormsNew}
\lean{HeckeRing.GL2.cuspFormsNew}
\uses{def:cuspFormsOld}
\leanok
The \emph{new subspace} $S_k^{\sharp}(N)$ is the orthogonal complement of the old subspace
$S_k^{\flat}(N)$ inside $S_k(\Gamma_1(N))$ with respect to the Petersson inner product,
\[
  S_k^{\sharp}(N) \;=\; \bigl(S_k^{\flat}(N)\bigr)^{\perp},
\]
following [DS, (5.19)]. Concretely, a cusp form lies in $S_k^{\sharp}(N)$ exactly when it is Petersson-orthogonal
to every oldform. The Petersson form is positive definite on cusp forms, so $S_k^{\sharp}(N)\cap S_k^{\flat}(N)=0$.
\end{definition}

\begin{definition}[Old part]\label{def:oldPart}
\lean{HeckeRing.GL2.oldPart}
\uses{def:cuspFormsOld, def:cuspFormsNew}
\leanok
Since $S_k(\Gamma_1(N))=S_k^{\flat}(N)\oplus S_k^{\sharp}(N)$ as an orthogonal direct sum, every cusp form
$f$ has a unique decomposition into an old and a new component. The \emph{old part} of $f$ is its image
under the projection of $S_k(\Gamma_1(N))$ onto $S_k^{\flat}(N)$ along $S_k^{\sharp}(N)$.
\end{definition}

\begin{definition}[New part]\label{def:newPart}
\lean{HeckeRing.GL2.newPart}
\uses{def:cuspFormsOld, def:cuspFormsNew}
\leanok
The \emph{new part} of a cusp form $f$ is its image under the projection of $S_k(\Gamma_1(N))$ onto the
new subspace $S_k^{\sharp}(N)$ along the old subspace $S_k^{\flat}(N)$. Together with the old part it
recovers $f$: one has $f=\operatorname{old}(f)+\operatorname{new}(f)$ with $\operatorname{old}(f)\in S_k^{\flat}(N)$
and $\operatorname{new}(f)\in S_k^{\sharp}(N)$.
\end{definition}

\begin{definition}[Newform]\label{def:newform}
\lean{HeckeRing.GL2.Newform}
\uses{def:eigenform, def:cuspFormsNew}
\leanok
A \emph{newform} (normalised newform) of level $\Gamma_1(N)$ and weight $k$ is an eigenform $f$ that lies
in the new subspace $S_k^{\sharp}(N)$ and is normalised so that its first Fourier coefficient is $a_1=1$.
By Atkin--Lehner uniqueness [DS, Thm~5.8.2] a newform is determined by its Hecke eigenvalues away from the
level; Theorem~\ref{thm:smo-clean} below is exactly this rigidity in the present formalisation.
\end{definition}

\section{Coefficients, leading terms, and the eigenform splitting}

\begin{theorem}[Miyake 4.5.15(1): $a_n=a_1\lambda_n$]\label{thm:coeff-eq-a1-lambda}
\lean{HeckeRing.GL2.Eigenform.coeff_eq_coeff_one_mul_eigenvalue}
\uses{def:eigenform, def:modFormCharSpace}
\leanok
Let $g$ be an eigenform lying in the Nebentypus space $S_k(N,\chi)$, and let $n$ be coprime to $N$. Then
the Fourier coefficients of $g$ satisfy
\[
  a_n(g) \;=\; a_1(g)\,\lambda_n(g),
\]
where $\lambda_n(g)$ is the classical Hecke eigenvalue of $g$ at $n$.
\end{theorem}

\begin{proof}
\uses{def:eigenform, def:modFormCharSpace}
\leanok
This is [Miy, Lemma~4.5.15(1)] in its un-normalised form. Apply the Hecke operator $T(n)$ to $g$. On the
one hand the eigenform relation gives $T(n)g=\lambda_n(g)\,g$, so the first Fourier coefficient of $T(n)g$
is $\lambda_n(g)\,a_1(g)$. On the other hand the explicit formula for the Fourier coefficients of $T(n)g$
expresses $a_m(T(n)g)$ as a sum over the common divisors of $m$ and $n$; taking $m=1$ and using $(n,N)=1$
the sum collapses to its single $d=1$ term, which is precisely $a_n(g)$. Equating the two expressions for
the first coefficient of $T(n)g$ yields $a_n(g)=a_1(g)\,\lambda_n(g)$.
\end{proof}

\begin{theorem}[Miyake 4.6.11: $a_1\neq 0$ for new eigenforms]\label{thm:a1-ne-zero}
\lean{HeckeRing.GL2.coeff_one_ne_zero_of_mem_cuspFormsNew_of_eigen_unconditional}
\uses{def:eigenform, def:cuspFormsNew, thm:coeff-eq-a1-lambda}
\leanok
Let $g$ be a nonzero eigenform lying in $S_k(N,\chi)\cap S_k^{\sharp}(N)$, that is, a nonzero common
eigenfunction of the $T(n)$, $(n,N)=1$, in the new subspace. Then its leading Fourier coefficient is
nonzero, $a_1(g)\neq 0$.
\end{theorem}

\begin{proof}
\uses{def:eigenform, def:cuspFormsNew, thm:coeff-eq-a1-lambda}
\leanok
This is [Miy, Lemma~4.6.11]. Argue by contraposition: suppose $a_1(g)=0$. By
Theorem~\ref{thm:coeff-eq-a1-lambda} every coprime coefficient then vanishes,
$a_n(g)=a_1(g)\,\lambda_n(g)=0$ for all $(n,N)=1$, and the zeroth coefficient vanishes because $g$ is
cuspidal. A form whose Fourier coefficients all vanish away from the level necessarily lies in the old
subspace (this is the content of [Miy, Thm~4.6.8] in the per-character formulation), so $g\in S_k^{\flat}(N)$.
But $g$ also lies in the new subspace $S_k^{\sharp}(N)$, and the Petersson form is positive definite, so
$S_k^{\flat}(N)\cap S_k^{\sharp}(N)=0$ forces $g=0$, contradicting the hypothesis. Hence $a_1(g)\neq 0$.
\end{proof}

\begin{theorem}[Eigenform old/new splitting]\label{thm:eigenform-decomp}
\lean{HeckeRing.GL2.exists_eigenform_decomposition_mem_cuspFormsNew}
\uses{def:eigenform, def:cuspFormsNew, def:oldPart, def:newPart}
\leanok
Let $g$ be a cusp form in the new subspace $S_k^{\sharp}(N)$ whose underlying modular form lies in the
Nebentypus space $S_k(N,\chi)$. Then $g$ is a finite sum of common Hecke eigenforms,
\[
  g \;=\; \sum_{i} h_i,
\]
each summand $h_i$ again lying in $S_k(N,\chi)\cap S_k^{\sharp}(N)$ and being a common eigenfunction of all
$T(n)$ with $(n,N)=1$.
\end{theorem}

\begin{proof}
\uses{def:eigenform, def:cuspFormsNew, def:oldPart, def:newPart}
\leanok
The new subspace is stable under every Hecke operator $T(n)$ with $(n,N)=1$ (DS Prop~5.6.2), and so is the
Nebentypus eigenspace $S_k(N,\chi)$. Since the operators $\{T(n)\}_{(n,N)=1}$ form a commuting family of
diagonalisable operators on the finite-dimensional space $S_k(N,\chi)$
(Theorem~\ref{thm:commring-hecke-algebra}), the simultaneous spectral decomposition of $S_k(N,\chi)$ into
joint eigenspaces restricts to the invariant subspace $S_k^{\sharp}(N)\cap S_k(N,\chi)$. Writing $g$ in
this joint eigenbasis exhibits it as a finite sum of common eigenforms, each of which is itself new and of
Nebentypus $\chi$.
\end{proof}

\section{The two halves of Theorem 4.6.12}

\begin{theorem}[New part is a multiple of $f$]\label{thm:newPart-eq-smul}
\lean{HeckeRing.GL2.newPart_eq_smul_of_shared_eigenvalues}
\uses{def:newform, def:newPart, thm:coeff-eq-a1-lambda, thm:a1-ne-zero, thm:eigenform-decomp}
\leanok
Let $f$ be a normalised newform in $S_k(N,\chi)$ and let $g^{(1)}$ be a common Hecke eigenfunction in the
new subspace lying in $S_k(N,\chi)$ that shares $f$'s eigenvalues $\lambda_n$ for all $(n,N)=1$ outside a
finite set $S$. Then
\[
  g^{(1)} \;=\; b_1\,f, \qquad b_1=a_1(g^{(1)}).
\]
\end{theorem}

\begin{proof}
\uses{def:newform, def:newPart, thm:coeff-eq-a1-lambda, thm:a1-ne-zero, thm:eigenform-decomp}
\leanok
This is the new-part half of [Miy, Thm~4.6.12]. If $g^{(1)}=0$ then $b_1=0$ and the identity is trivial,
so assume $g^{(1)}\neq 0$; then $b_1\neq 0$ by Theorem~\ref{thm:a1-ne-zero}. Sharing eigenvalues off the
finite set $S$ in fact forces equality of \emph{all} coprime eigenvalues (a nonzero coprime eigenvalue and
the multiplicativity of the eigenvalue system propagate the agreement to every $(n,N)=1$). Consider
$h=b_1^{-1}g^{(1)}-f$. For $(n,N)=1$ its $n$-th coefficient is
$b_1^{-1}a_n(g^{(1)})-a_n(f)=\lambda_n-\lambda_n=0$, using Theorem~\ref{thm:coeff-eq-a1-lambda} together
with the normalisation $a_1(f)=1$ which gives $a_n(f)=\lambda_n$. A form with vanishing coprime
coefficients lies in the old subspace by [Miy, Thm~4.6.8], so $h\in S_k^{\flat}(N)$; but $h$ is a
difference of new forms, hence $h\in S_k^{\sharp}(N)$. Orthogonality of the new and old subspaces forces
$h=0$, i.e.\ $b_1^{-1}g^{(1)}=f$, which is the claim.
\end{proof}

\begin{theorem}[Old part vanishes]\label{thm:oldPart-eq-zero}
\lean{HeckeRing.GL2.oldPart_eq_zero_of_shared_eigenvalues}
\uses{def:newform, def:oldPart, thm:a1-ne-zero, thm:eigenform-decomp}
\leanok
Let $f$ be a nonzero newform in $S_k(N,\chi)$ and let $g^{(2)}$ be a common Hecke eigenfunction in the
($\chi$-refined) old subspace, lying in $S_k(N,\chi)$, that shares $f$'s eigenvalues $\lambda_n$ for all
$(n,N)=1$ outside a finite set $S$. Then $g^{(2)}=0$.
\end{theorem}

\begin{proof}
\uses{def:newform, def:oldPart, thm:a1-ne-zero, thm:eigenform-decomp}
\leanok
This is steps (i)--(ii) of [Miy, Thm~4.6.12]. Suppose $g^{(2)}\neq 0$. By the description of the old
subspace [Miy, Lemma~4.6.9] one may write $g^{(2)}$ as a sum of level-raises $h_v(\ell_v z)$ of newforms
$h_v$ at proper divisors $M_v$ of $N$; by Theorem~\ref{thm:eigenform-decomp} the $h_v$ may be taken to be
common eigenforms, and by [Miy, Lemma~4.6.2] their level-raises remain eigenforms with the same eigenvalues.
Eigenfunctions with distinct eigenvalues are Petersson-orthogonal, hence linearly independent, so the
components whose eigenvalues differ from those of $f$ contribute nothing; removing them leaves a nonzero
new eigenform $h$ at a proper divisor $M$ whose eigenvalues match $f$'s off $S$. By
Theorem~\ref{thm:a1-ne-zero} its leading coefficient $c_1'$ is nonzero. The difference $h-c_1'f$ has
vanishing coprime coefficients (Theorem~\ref{thm:coeff-eq-a1-lambda} and $a_1(f)=1$), so it lies in the old
subspace by [Miy, Thm~4.6.8]; and $h$ itself, being a form essentially of level $M\neq N$, also lies in the
old subspace. Therefore
\[
  f \;=\; c_1'^{-1}h - c_1'^{-1}(h-c_1'f)
\]
exhibits $f$ as an element of the old subspace. This contradicts $f$ being a nonzero newform (the new and
old subspaces meet only in $0$). Hence $g^{(2)}=0$.
\end{proof}

\section{Strong Multiplicity One}

\begin{theorem}[Strong Multiplicity One, Miyake 4.6.12]\label{thm:smo-constMul}
\lean{HeckeRing.GL2.strongMultiplicityOne_constMul}
\uses{def:newform, def:eigenform, def:oldPart, def:newPart, thm:newPart-eq-smul, thm:oldPart-eq-zero}
\leanok
Let $f$ be a normalised newform in $S_k(N,\chi)$, and let $g$ be any cusp-form Hecke eigenfunction in
$S_k(N,\chi)$. If $f$ and $g$ are common eigenfunctions of $T(n)$ with the same eigenvalue $\lambda_n$ for
every $n$ with $(n,N)=1$ outside a finite set $S$, then $g$ is a constant multiple of $f$:
\[
  g \;=\; c\,f \qquad\text{for some } c\in\mathbb{C}.
\]
\end{theorem}

\begin{proof}
\uses{def:newform, def:eigenform, def:oldPart, def:newPart, thm:newPart-eq-smul, thm:oldPart-eq-zero}
\leanok
This is [Miy, Thm~4.6.12]. Normalise $a_1(f)=1$, which is legitimate since $a_1(f)\neq 0$ by
[Miy, Lemma~4.6.11], and split $g$ along the orthogonal new/old decomposition,
$g=g^{(1)}+g^{(2)}$ with $g^{(1)}\in S_k^{\sharp}(N)$ and $g^{(2)}\in S_k^{\flat}(N)$. Because both the new
subspace and the old subspace are stable under the Hecke operators (DS Prop~5.6.2 / [Miy, Lemma~4.6.10]),
each of $g^{(1)}$ and $g^{(2)}$ is again a common eigenfunction of $T(n)$ sharing $f$'s eigenvalue
$\lambda_n$ for $(n,N)=1$ outside $S$.

For the new part, Theorem~\ref{thm:newPart-eq-smul} gives $g^{(1)}=b_1\,f$ where $b_1=a_1(g^{(1)})$: the
difference $b_1^{-1}g^{(1)}-f$ has vanishing coprime coefficients by [Miy, Lemma~4.5.15(1)] and therefore
lies in the old subspace by [Miy, Thm~4.6.8], while being new, hence is $0$. For the old part,
Theorem~\ref{thm:oldPart-eq-zero} gives $g^{(2)}=0$: assuming otherwise, one descends $g^{(2)}$ to a
nonzero new eigenform $h$ at a proper divisor sharing $f$'s eigenvalues ([Miy, Lemma~4.6.9],
[Miy, Lemma~4.6.2], and linear independence of distinct-eigenvalue eigenforms), and then
$h-c_1'f\in S_k^{\flat}(N)$ together with $h\in S_k^{\flat}(N)$ forces $f\in S_k^{\flat}(N)$, contradicting
that $f$ is a nonzero newform. Combining the two halves, $g=g^{(1)}+g^{(2)}=b_1\,f$, so $g=c\,f$ with
$c=b_1=a_1(g)$.
\end{proof}

\begin{theorem}[Multiplicity One for newforms, DS 5.8.2]\label{thm:smo-clean}
\lean{HeckeRing.GL2.strongMultiplicityOne_axiom_clean_unconditional}
\uses{thm:smo-constMul, def:newform}
\leanok
Let $f$ and $g$ be normalised newforms in the same Nebentypus space $S_k(N,\chi)$. If $f$ and $g$ have the
same Hecke eigenvalue $\lambda_n$ for every $n$ with $(n,N)=1$ outside a finite set $S$, then $f=g$.
\end{theorem}

\begin{proof}
\uses{thm:smo-constMul, def:newform}
\leanok
This is [DS, Thm~5.8.2] and follows immediately from the headline Theorem~\ref{thm:smo-constMul}. Applying
it to the newform $f$ and the eigenfunction $g$ produces a constant $c$ with $g=c\,f$. Comparing first
Fourier coefficients and using that both forms are normalised, $1=a_1(g)=c\,a_1(f)=c$, so $c=1$ and hence
$g=f$.
\end{proof}

